\documentclass{article} % say
\usepackage{amsmath}
\usepackage{gensymb}
\usepackage{hyperref}
\usepackage{tikz}
\usepackage{tkz-euclide}
\usepackage{amsfonts}
\usetikzlibrary {babel,angles,quotes,positioning}
\usepackage[utf8]{inputenc}
\usepackage[russian]{babel}
\begin{document}

\textbf{Задача}: Докажите, что $\cos 1\degree$ иррационален.

\textbf{Подсказка}: Докажите равенство:

\begin{equation}
\cos(nx) = 2\cos(x) \cos\left( (n-1) x \right) - \cos\left( (n-2) x \right) \label{eq:hint}
\end{equation}

\textbf{Решение}: Воспользуемся формулами косинуса суммы и разницы:

\begin{align*}
  \cos(\alpha + \beta) &= \cos\alpha \cos\beta - \sin\alpha \sin\beta \\
  \cos(\alpha - \beta) &= \cos\alpha \cos\beta + \sin\alpha \sin\beta
\end{align*}

Примем $\alpha = (n-1)x$, $\beta=x$, тогда:

\begin{align}
  \cos(nx) &= \cos((n-1)x) \cos x - \sin((n-1)x) \sin x \label{eq:sum} \\
  \cos((n-2)x) &= \cos((n-1)x) \cos x + \sin((n-1)x) \sin x \label{eq:diff}
\end{align}

При сложении уравнений \ref{eq:sum} и \ref{eq:diff} получаем уравнение \ref{eq:hint}.

Теперь положим, что $\cos 1\degree$ рационален. Покажем, что тогда $\cos k\degree$, где $k \in \mathbb{N}$, также рационален. Отсюда придём к противоречию, например при $k=30$.

Действительно, пусть $\cos x$ рационален, докажем, что $\cos nx$ рационален при любом $n \in \mathbb{N}$:

\begin{itemize}
\item при $n=1$ утверждение тривиально;
\item при $n=2$: $\cos(2x) = 2\cos^2x - 1$, т.~е. рационален;
\item при $n>2$ воспользуемся индукцией: по уравнению \ref{eq:hint} $\cos(nx)$~--- это рациональное выражение над рациональными числами, т.~е. рациональное число.
\end{itemize}

\textbf{Источник}: \url{https://t.me/obznam/1409}

\end{document}


